\documentclass[10pt,a4paper]{amsart}
\usepackage[margin=19mm]{geometry}
\usepackage{amsmath,amssymb,mathtools}
\usepackage[scr=rsfs,cal=euler]{mathalfa}
\usepackage{xfrac}
\usepackage[unicode,hidelinks,bookmarks=false]{hyperref}
\newcommand{\ie}{i.e.\ }
\newcommand{\cf}{cf.\ }
\newcommand{\resp}{respectively\ }
\newcommand{\Subsection}{\subsection}
\providecommand{\nobreakdash}{\nobreak\hbox{-}}
\setlength{\emergencystretch}{2em}
\multlinegap0pt
\usepackage{mdframed}
\usepackage{mdframed}
\usepackage{mdframed}
\usepackage{mdframed}
\usepackage{braket}
\usepackage{amssymb}
\usepackage{enumerate}
\usepackage{float}
\usepackage{mathtools}
\usepackage{mdframed}
\usepackage{xcolor}
\usepackage{algorithm}\usepackage{algpseudocode}
\newtheorem{lemma}{Lemma}
\title{Large sets of mutually orthogonal quantum Latin squares}
\author{Simeon Ball, Robin Simoens}
\date{Source: arXiv:2607.12933v2 (CC BY 4.0)}
\begin{document}
\maketitle

\noindent\emph{Source and licence.} Large sets of mutually orthogonal quantum Latin squares. Version arXiv:2607.12933v2 (CC BY 4.0), distributed by arXiv under the Creative Commons Attribution 4.0 International licence, http://creativecommons.org/licenses/by/4.0/, as recorded in the arXiv licence metadata for this exact version and shown on the abstract page https://arxiv.org/abs/2607.12933v2. Full licence notice: \texttt{licenses/arxiv-2607.12933-CC-BY-4.0.txt}.\par\smallskip

\noindent\textit{Editorial scope.} Counted unit: the article's lemma lem:classical (source0.tex lines 341--343). The article's complete original proof of it is delivered with it (source0.tex lines 344--358, one \texttt{proof} environment of 15 lines). No mathematics is written, repaired or re-derived: the statement and the proof are the article's own lines, in the article's own notation, and no retained line is substituted or re-flowed.  No further source-local context is needed: every cross-reference of every retained line resolves inside the two delivered spans.  Nothing was omitted from any retained span, so the package carries no omission record.  No retained line contains a citation, so the wrapper carries no bibliography and no bibliography entry.  No retained line contains a figure environment, an image-inclusion command or drawing code, so no image is delivered and the package declares no asset dependency.  Editorial formatting only, in addition: an AMS \texttt{amsart} preamble that restates the article's own declarations and macros used by the retained lines, namely the environments \texttt{lemma} on separate counters, together with the standard packages those lines need.  The retained environments print in this document as Lemma 1 in order of appearance, whereas the article prints the same items with the numbering of its own full document.  The complete native source of the article as distributed in the arXiv e-print for this exact version is bundled unchanged as \texttt{sources/205153/source0.tex}.
\begin{lemma}\label{lem:classical}
    Let \(f:R\to R\) be a permutation. The quantum Latin square \(\Psi^{(f)}\) is classical if and only if \(f\) is the sum of an additive map and a constant.
\end{lemma}

\begin{proof}
    Replacing \(f\) by \(f-f(0)\) changes every entry \(\psi_{ij}^{(f)}\) only by the phase factor \(\chi(if(0))\), so assume that \(f(0)=0\).
    Since every row of a quantum Latin square of order \(n\) already contains \(n\) distinct entries, \(\Psi^{(f)}\) is classical if and only if its total cardinality (the number of distinct entries up to phase) is \(n\).
    %The first row of \(\Psi^{(f)}\) contains the \(n\) different entries \(\frac1{\sqrt{n}}\sum_{x\in R}\chi(jx)\ket{x}\).

    First assume that \(\Psi^{(f)}\) is classical. The first row and column must have the same entries up to phase, so for every \(i\in R\), there exists a phase factor \(c\in\mathbb{T}\) and a column index \(j\in R\) such that
    \[\frac1{\sqrt{n}}\sum_{x\in R}\chi(if(x))\ket{x}=\frac{c}{\sqrt{n}}\sum_{x\in R}\chi(jx)\ket{x}.\]
    Comparing coefficients gives that, for every \(i\in R\), there are \(c\in\mathbb{T}\) and \(j\in R\) such that
    \[\chi(if(x))=c\chi(jx)\]
    for all \(x\in R\).
    Filling in \(x=0\), we get that \(c=1\) because \(f(0)=0\) and \(\chi(0)=1\). 
    Since \(x\mapsto \chi(jx)\) is an additive character, the map \(x\mapsto\chi(if(x))\) is additive character. That is, \(\chi(if(x+y))=\chi(if(x))\cdot\chi(if(y))\) for all \(i,x,y\in R\). By additivity of \(\chi\), we get \(\chi(i(f(x+y)-f(x)-f(y)))=1\) for all \(i,x,y\in R\). But \(\chi\) is a generating character, so \(f(x+y)=f(x)+f(y)\) for all \(x,y\in R\). In other words, \(f\) is additive.
    
    Conversely, if \(f\) is additive, then for every \(i\in R\), the map \(x\mapsto\chi(if(x))\) is an additive character, so by the defining property of a generating character, there exists a unique element \(a\in R\) such that \(\chi(if(x))=\chi(ax)\) for \(x\in R\). Thus, \(\psi_{ij}^{(f)}=\frac1{\sqrt{n}}\sum_{x\in R}\chi((a+j)x)\ket{x}\) is equal to the first-row entry on position \((0,a+j)\). We conclude that \(\Psi^{(f)}\) has cardinality \(n\) and is therefore classical.
\end{proof}

\end{document}
